\chapter{Results}
\label{sec:results}

This chapter collects the validation and production measurements. The
harmonic backbone is checked first (\cref{sec:res_ballistic}), then the
force-constant inputs (\cref{sec:res_fc}), the anharmonic transport on
finite silicon nanowires (\cref{sec:res_nw}), the transversely-periodic
$q_\perp$-resolved anharmonic problem and the silicon thin-film benchmark
(\cref{sec:res_qfilm}), the still-open question of the absolute
three-phonon self-energy prefactor (\cref{sec:res_prefactor}), the cost
of the numerical approximations (\cref{sec:res_approx}), a Si/Ge
heterostructure (\cref{sec:res_hetero}), and the distributed-solver and
GPU-memory scaling of the implementation
(\cref{sec:res_scaling}). Implementation mechanisms referenced here are
described in \cref{ssec:phph_distribution,ssec:phph_qdist}; the dense
reference solver used for the standalone studies is documented in
\cref{app:convergence_diagnosis}.

% ======================================================================
\section{Ballistic phonon NEGF validation}
\label{sec:res_ballistic}
% ======================================================================

The harmonic transport backbone is validated against an independent
Sancho--Rubio decimation plus Caroli reference on identical harmonic
inputs, transverse mesh, and frequency grid. For homogeneous bulk
silicon the two transmissions agree, confirming the phonopy block
extraction, the block-tridiagonal assembly, and the open-boundary
treatment.

As a heterogeneous test we model a ballistic Si/Ge interface in the
mass-mismatch approximation: bulk-Si force constants throughout, with
the Si/Ge distinction entering only through the atomic masses
($m_\mathrm{Si}=28.09$, $m_\mathrm{Ge}=72.64$~amu). At $\Gamma$ the Si
optical maximum is $15.55$~THz and the Ge mass-clone gives $9.67$~THz,
consistent with $f_\mathrm{Ge}=f_\mathrm{Si}\sqrt{m_\mathrm{Si}/m_\mathrm{Ge}}$.
On a $20\times20$ transverse mesh the interfacial conductance at
$300$~K is $G=196$~MW\,m$^{-2}$K$^{-1}$, within $\sim7\%$ of the
literature values ($210$, Latour 2017; $209$, Tian 2012); the
remaining gap is consistent with the coarse transverse mesh and the
minimal FC2 calculation.

% ======================================================================
\section{Force constants: methods, quality, and compression}
\label{sec:res_fc}
% ======================================================================

\paragraph{Method cross-check.} Bulk-silicon harmonic force constants
computed by density-functional perturbation theory (QE \texttt{ph.x}
$\to$ \texttt{q2r.x}) and by finite displacement (phono3py + symfc) on
the same primitive cell agree to $1.6\%$ on the zone-centre optical
frequency ($15.12$ vs $15.37$~THz; \cref{fig:res_fcmethod}). The two
\emph{ab-initio} routes are therefore mutually consistent at the
harmonic level. (DFPT third-order constants via \texttt{d3q} require a
norm-conserving pseudopotential and were not used for transport here.)

\begin{figure}[htb]
  \centering
  \includegraphics[width=0.6\linewidth]{transport_sweeps/fc_method_dispersion_si.pdf}
  \caption{Bulk-silicon phonon dispersion from DFPT versus
  finite-displacement force constants. The two \emph{ab-initio} methods
  agree to $1.6\%$ on the $\Gamma$ optical frequency.}
  \label{fig:res_fcmethod}
\end{figure}

\paragraph{Bulk-$\kappa$ anchor.} The finite-displacement FC3, fed to a
phono3py RTA-BTE calculation (natural isotopes), gives bulk-silicon
$\kappa(300\,\mathrm{K})=110.7$~W\,m$^{-1}$K$^{-1}$ on an $11^3$ mesh,
rising to $115.2$ on $19^3$. This is $\sim20$--$30\%$ below the
experimental $\sim150$, the known underestimate of a $2\times2\times2$
FC3 supercell (cutoff $5.5$~\AA); a larger supercell would close most of
the gap. This bulk $\kappa$ is the reference against which the
anharmonic self-energy prefactor should ultimately be benchmarked
(\cref{sec:res_prefactor}).

\paragraph{FC3 compression.} The decomposition machinery of
\cref{sec:fc3_compression} was applied to the $3\times3\times3$ silicon
FC3. The transport observable converges far faster than the Frobenius
residual: targeting a $1\%$ relative conductance error is reached by an
INDSCAL fit at rank $8$ ($117\times$ compression) while its Frobenius
residual is still $\sim17\%$, and the spectral heat-current curves are
visually indistinguishable from the dense reference by rank $16$. The
practical lesson is that the relevant fitting target for phonon NEGF is
a transport residual, not the Frobenius norm the fitters minimise.

% ======================================================================
\section{Anharmonic transport in silicon nanowires}
\label{sec:res_nw}
% ======================================================================

We study two hydrogen-passivated [100] silicon nanowires: \texttt{d5a}
($\mathrm{Si}_9\mathrm{H}_{12}$ per cell, $63$ DOF) and the wider
\texttt{d11a} ($135$ DOF). Their harmonic transport is clean
(\cref{fig:res_ballistic_curves}): the ballistic conductance rises with
temperature (thermal activation, saturating) and \texttt{d11a} carries
$\sim1.5\times$ the conductance of \texttt{d5a} (more channels). The
apparent decrease of $G$ with device length at finite broadening is the
$\eta$-induced coherence attenuation ($\ell\sim v_g/\eta_w$), not
diffusive resistance; the true Landauer plateau requires $\eta\to0$.

\begin{figure}[htb]
  \centering
  \includegraphics[width=0.48\linewidth]{transport_sweeps/ballistic_vs_T_d5_d11.pdf}\hfill
  \includegraphics[width=0.48\linewidth]{transport_sweeps/ballistic_vs_length_d5_d11.pdf}
  \caption{Ballistic conductance of the \texttt{d5a} and \texttt{d11a}
  silicon nanowires. Left: $G(T)$; right: $G$ versus device length (the
  decrease is finite-$\eta$ coherence attenuation).}
  \label{fig:res_ballistic_curves}
\end{figure}

\paragraph{Mode-resolved anharmonic physics.} A phono3py analysis of the
bulk-silicon FC3 (\cref{fig:res_modephysics}) gives the context for the
nanowire results: phonon lifetimes have a median of $4.2$~ps with a
$\tau\!\propto\!\omega^{-2}$ envelope; $50\%$ of $\kappa$ is carried by
phonons with mean free path $>134$~nm and $93\%$ by acoustic modes;
Umklapp accounts for $0.44$ of the scattering at $300$~K; and the mode
Gr\"uneisen parameters have mean $|\gamma|=1.04$, matching the
experimental Si value and independently validating the FC3
\emph{anharmonicity strength}. Crucially, only $1.3\%$ of bulk
$\kappa$ comes from phonons with mean free path below $10$~nm. The
$\sim1$~nm nanowire cores therefore sit far below the bulk
phonon-MFP spectrum: they are \emph{boundary-scattering limited} and
retain almost none of the bulk anharmonic-limited $\kappa$. This is
why, at a single transport cell, the anharmonic correction to the
nanowire conductance is a small perturbation
($G_\mathrm{anh}/G_\mathrm{ball}\approx1$), and it is consistent across
the BTE-RTA and NEGF pictures.

\begin{figure}[htb]
  \centering
  \includegraphics[width=0.48\linewidth]{transport_sweeps/phph_physics_si.pdf}\hfill
  \includegraphics[width=0.48\linewidth]{transport_sweeps/phph_NU_gruneisen_si.pdf}
  \caption{Bulk-silicon mode-resolved anharmonic physics (phono3py,
  FD FC3). Left: lifetimes, $\kappa$-accumulation versus mean free path,
  and the acoustic-mode $\kappa$ fraction. Right: Normal/Umklapp split
  and the mode Gr\"uneisen distribution (mean $|\gamma|=1.04$).}
  \label{fig:res_modephysics}
\end{figure}

At physical broadening and moderate vertex strength the SCBA converges
cleanly and gives the physical $G_\mathrm{anh}<G_\mathrm{ball}$; the
soft-mode (\texttt{d5a} $0.027$~THz twist) and strong-coupling
behaviour, and the diagnosis that the original $\eta/\lambda$ sweep was
run at an unphysically large $\eta$, are documented in
\cref{app:convergence_diagnosis}.

% ======================================================================
\section{q-resolved anharmonic transport: silicon thin film}
\label{sec:res_qfilm}
% ======================================================================

For a transversely-periodic system the three-phonon self-energy couples
the transverse momenta by crystal-momentum conservation
(\cref{eq:sigma_q_dist}). As a quantitative benchmark we model
cross-plane heat transport through a silicon thin film following
Guo, Bescond and Zhang~\cite{guo_anharmonic_negf_2020}: the film is
bulk silicon truncated to $N$ layers along the transport direction with
in-plane two-dimensional periodicity, so it is exactly the
transverse-$q_\perp$ problem and reuses the \emph{bulk} silicon force
constants. The driver reduces to the $\Gamma$-only finite-device
transmission at a $1\times1$ mesh (checked automatically). Conductance
is reported per unit area, $G=J/(A_\perp\Delta T)$, and the cross-plane
conductivity of a film of thickness $L$ is $\kappa(L)=G\,L$.

\paragraph{Ballistic baseline.} The transverse mesh and broadening must
both be converged: at a fixed eight-layer thickness the ballistic
conductance converges for $8\times8$ or finer
($G_\mathrm{ball}=970/928/921/920$~MW\,m$^{-2}$K$^{-1}$ at
$4\times4/8\times8/12\times12/16\times16$), and a small $\eta$ is needed
to avoid the coherence attenuation of \cref{sec:res_nw} that otherwise
crushes the conductance. With these settings $G_\mathrm{ball}$ is nearly
thickness-independent ($\approx10^{3}$~MW\,m$^{-2}$K$^{-1}$) and the
ballistic $\kappa=G_\mathrm{ball}L$ rises essentially linearly with $L$
(\cref{fig:res_sifilm}, right) -- the non-intensive ballistic regime.
This ballistic conductance matches the scale of
Guo~\emph{et al.} to within a few percent.

\paragraph{Anharmonic conductance.} \Cref{fig:res_sifilm} (left)
compares our cross-plane conductance to the anharmonic-NEGF values of
Guo~\emph{et al.} at $3$ and $5$ unit cells. With the self-energy as
implemented (\emph{native} prefactor, \cref{sec:res_prefactor}) the
three-phonon scattering reduces the $3$-unit-cell conductance by
$\sim45\%$ ($G_\mathrm{anh}\approx570$~MW\,m$^{-2}$K$^{-1}$), markedly
more than the reference's $\sim10\%$
($G_\mathrm{anh}=939.7$). \emph{We have not resolved this
discrepancy}: it is either a transport-setup effect ($\eta$, mesh,
thin-device artifacts), the reference's additional approximations, the
difference in exchange--correlation functional of the underlying force
constants, or an absolute-prefactor error in the self-energy. The
second (green) curve in \cref{fig:res_sifilm} shows the effect of
scaling the vertex by one half (so $\Sigma\to\Sigma/4$); it moves the
conductance toward the reference ($831$ and $704$) but does not match
it, and we do \emph{not} claim it as a correction. The prefactor is
discussed honestly in \cref{sec:res_prefactor}.

\begin{figure}[htb]
  \centering
  \includegraphics[width=0.92\linewidth]{transport_sweeps/si_film_conductance.pdf}
  \caption{Cross-plane heat transport through a silicon thin film at
  $300$~K with bulk-silicon force constants (transverse mesh $8\times8$,
  $121$ frequencies). \emph{Left}: per-area conductance versus
  thickness -- our ballistic and anharmonic NEGF against
  Guo~\emph{et al.}~\cite{guo_anharmonic_negf_2020}. The native
  self-energy over-scatters relative to the reference; the
  vertex-scaled curve illustrates the unresolved absolute prefactor
  (\cref{sec:res_prefactor}), and is not an established correction.
  \emph{Right}: the ballistic $\kappa=G_\mathrm{ball}L$ rises linearly
  with thickness (ballistic regime).}
  \label{fig:res_sifilm}
\end{figure}

% ======================================================================
\section{The three-phonon self-energy prefactor (open)}
\label{sec:res_prefactor}
% ======================================================================

The absolute scale of the anharmonic self-energy is, at the time of
writing, \textbf{not settled}, and we report it as an open issue rather
than a resolved one. The implemented kernel is the standard sunset
diagram (\cref{eq:sigma_guo}), which reduces analytically to the
correct Fermi-golden-rule three-phonon linewidth
(\cref{eq:fgr}). Against this, the silicon thin film over-scatters
relative to Guo~\emph{et al.} (\cref{sec:res_qfilm}), which would be
reconciled by a factor in the range $\sim2$--$4$ on the self-energy
(Guo~\emph{et al.} note a ``factor of four'' over-count of the repeated
internal pairings in the earlier Luisier convention). However, our
attempt to pin this factor independently against phono3py's
golden-rule linewidths for bulk silicon on the same FC3 is dominated by
the equilibrium-Green's-function normalisation: the broadening-free,
units-matched ratio
$\int(-\mathrm{Im}\,\Sigma^{\mathrm{NEGF}}_s)\,d\omega/\int 2\omega_s\,\gamma^{\mathrm{p3p}}_s\,d\omega$
comes out mode-independent and close to $(2\pi)^2$, i.e.\ a
normalisation/unit artifact of the standalone reconstruction rather than
a clean physical factor. The check therefore does not determine the
prefactor, and we have reverted to the unscaled (native) self-energy as
the default. The definitive, convention-free test -- bulk-silicon
$\kappa$ computed from the code's \emph{own} SCBA Green's functions and
compared to the phono3py RTA value of \cref{sec:res_fc} on the identical
FC3 -- is identified but not yet completed. All anharmonic conductances
in this chapter are quoted at the native prefactor and should be read
with this caveat.

% ======================================================================
\section{Full versus approximated self-energy}
\label{sec:res_approx}
% ======================================================================

The $q_\perp$-resolved transport driver computes, by construction, only
the slab-diagonal self-energy (the reference's ``diagonal-block''
approximation). We implemented the \emph{full} off-diagonal multi-slab
self-energy $\Sigma_{IJ}(q_\perp,\omega)$ -- which reduces exactly
(relative error $0$) to the $\Gamma$-only multi-slab self-energy at a
$1\times1$ mesh -- to measure each numerical approximation against the
unapproximated result. These ratios are prefactor-independent. For a
three-slab silicon film (first-Born step from the ballistic Green's
function): the off-diagonal self-energy blocks carry $47\%$ of the
Frobenius weight yet change the conductance by only $-1.8\%$, so the
diagonal-block approximation is well justified despite the blocks being
individually large; restricting the third-order vertex to the
nearest-neighbour slab shell changes the result by $<0.1\%$ (the
long-range FC3 is negligible, echoing the FC3-quality analysis); but
truncating the inter-slab \emph{Green's function} range gives an
unphysical (negative) conductance -- the off-diagonal $G$ coherence is
essential and must not be cut. Of the reference's numerical
approximations, the two it makes are benign at the few-percent level
here, while the Green's-function-range truncation, which it does
\emph{not} make, would be catastrophic.

% ======================================================================
\section{Silicon/germanium heterostructure}
\label{sec:res_hetero}
% ======================================================================

A Ge-layer barrier embedded in a silicon film (Si force constants,
Ge mass; Si contacts) is built with a per-slab mass profile threaded
through the $q_\perp$-resolved driver (an all-Si profile reproduces the
homogeneous driver to machine precision). The mass-mismatch barrier
raises the \emph{ballistic} thermal resistance by $+340$--$370\%$
(harmonic interface reflection). The phonon-phonon enhancement of the
resistance is present in the pure silicon film -- $+32\%$ at $2.3$~nm
growing to $+51\%$ at $5.4$~nm (the anharmonic resistance accumulates
with thickness, the diffusive crossover) -- but is suppressed
($\lesssim1\%$) in the barrier-dominated heterostructure, because the
strong mass barrier filters transmission to low-frequency acoustic
phonons that are anharmonically long-lived. (The magnitude of the
anharmonic enhancement inherits the prefactor caveat of
\cref{sec:res_prefactor}; the barrier reflection and the qualitative
filtering picture do not.)

% ======================================================================
\section{Distributed solver and self-energy; GPU memory}
\label{sec:res_scaling}
% ======================================================================

\paragraph{Dyson solve.} The block-tridiagonal recursive Green's
function (RGF) solve is linear in the number of transport cells while
the dense reference is cubic ($57\times$ apart at $32$ blocks, agreeing
to $\sim10^{-13}$; \cref{fig:res_solver}, left). Distributing the energy
points across MPI ranks gives $5.3\times$ at eight ranks
(\cref{fig:res_solver}, right) -- the Dyson solve is
energy-embarrassingly-parallel.

\begin{figure}[htb]
  \centering
  \includegraphics[width=0.48\linewidth]{transport_sweeps/solver_scaling.pdf}\hfill
  \includegraphics[width=0.48\linewidth]{transport_sweeps/phph_scaling.pdf}
  \caption{Left: solver scaling -- RGF (linear) versus dense
  selected-inversion (cubic), and the energy-parallel speed-up
  ($5.3\times$ at eight ranks). Right: distributed scaling of the
  production three-phonon self-energy -- the spatial-block
  $(I,J)$ axis is near-ideal while the energy axis is flat because the
  energy convolution is replicated (\cref{ssec:phph_distribution}).}
  \label{fig:res_solver}
\end{figure}

\paragraph{Self-energy.} The production phonon-phonon self-energy
parallelises near-ideally over the spatial-block axis ($1.0/1.52/2.37/5.36\times$
at $1/2/3/6$ ranks) but is flat over the energy axis: the bubble is an
energy convolution that couples all frequencies, so each stack rank
recomputes the full-frequency FFT and writes only its slice
(\cref{fig:res_solver}, right).

\paragraph{Transverse-momentum axis.} The dedicated $q$-communicator of
\cref{ssec:phph_qdist} -- distributing the external $q_\perp$ and
all-gathering the internal-momentum Green's functions -- scales
near-ideally: $1.0/1.87/3.58/6.58\times$ at $1/2/4/8$ ranks
($82\%$ efficiency), with the self-energy bit-identical across rank
counts (\cref{fig:res_qscaling}, left). This is the scalable axis for
the periodic anharmonic problem, in contrast to the replicated energy
axis.

\paragraph{GPU memory.} For the periodic problem the binding constraint
is memory: the dense Fourier vertex $\tilde\Phi(q_\perp,q_\perp')$ scales
as $\cO(N_q^2\,\ndof^3)$ (16~GB already at an $8\times8$ mesh and
wire-scale cross section). Streaming it from the gathering matrices and
the real-space FC3 (\cref{ssec:phph_qdist}) drops the peak to
$\cO(N_q\,\ndof^2)$ (\cref{fig:res_qscaling}, right), bit-identical to
the dense path, and composes with the $q$-communicator at no extra
communication.

\begin{figure}[htb]
  \centering
  \includegraphics[width=0.48\linewidth]{transport_sweeps/phph_q_scaling.pdf}\hfill
  \includegraphics[width=0.48\linewidth]{transport_sweeps/phph_memory.pdf}
  \caption{Left: distributed scaling of the $q_\perp$-resolved
  three-phonon self-energy over the $q$-communicator (near-ideal;
  self-energy bit-identical across ranks). Right: peak memory of the
  momentum-space vertex -- dense $\cO(N_q^2\ndof^3)$ versus streamed
  $\cO(N_q\ndof^2)$ (the star is a measured value).}
  \label{fig:res_qscaling}
\end{figure}

% ======================================================================
\section{Cutoff hierarchy and transport sensitivity}
\label{sec:res_cutoffs}
% ======================================================================

A self-energy-level cutoff sweep on the \texttt{d5a} and \texttt{d11a}
nanowires (\cref{fig:res_cutoffs}) quantifies the truncations: the FC3
magnitude threshold converges quickly for \texttt{d5a}
($10^{-3}\to0.7\%$ relative change in $\Sigma^<$) but the wider
\texttt{d11a} needs a tighter threshold (it has more small-but-relevant
vertex elements), whereas a diagonal-$G$ approximation inside the bubble
is catastrophic ($5\times$ for \texttt{d5a}, $30\times$ for
\texttt{d11a}). The dominant cutoff is the Green's-function range, not
the vertex magnitude -- consistent with the q-resolved finding of
\cref{sec:res_approx} that the inter-slab $G$ coherence must be kept.

\begin{figure}[htb]
  \centering
  \includegraphics[width=0.6\linewidth]{transport_sweeps/cutoff_sse_d5_d11.pdf}
  \caption{Self-energy sensitivity to the FC3 magnitude cutoff and to a
  diagonal-$G$ approximation for the \texttt{d5a} and \texttt{d11a}
  nanowires. The off-diagonal Green's-function content dominates.}
  \label{fig:res_cutoffs}
\end{figure}
